%! TeX program = lualatex
\documentclass[a4paper, 14pt]{article}
\usepackage{fontspec}
\usepackage[14pt]{extsizes}
\usepackage[utf8x]{inputenc}
\usepackage[russian]{babel}
\usepackage{setspace}
\usepackage{indentfirst}

\setlength{\parindent}{12.5mm}
\usepackage{multirow}
\usepackage{tabularx}
\usepackage{pdflscape}
\usepackage{geometry}
\geometry{top=2cm}
\geometry{right=1cm}
\geometry{left=3cm}
\geometry{bottom=2cm}
\setmainfont{Times New Roman}
\usepackage{amsmath}
\usepackage{amsfonts}
\usepackage{graphicx}
\usepackage{xcolor}
\usepackage{xlop}
\usepackage{amsmath}
\date{}

\usepackage{array}
\usepackage{makecell}
\usepackage{etoolbox}
\AtBeginEnvironment{tabular}{\noindent}

\usepackage{xfp}
\usepackage{siunitx}
\usepackage{etoolbox}
\usepackage{subcaption}
\usepackage{float}
\usepackage[hidelinks]{hyperref}

\usepackage{etuphyspreparation}

\begin{document}

\labpreparation{number=0,labname=Название работы,myname=Иванов Иван Иванович,group=4308,faculty=Ф\ldots,professor=Григорьев Григорий Григорьевич,year=2024}{
  Цель работы
} {
  Краткое описание
  \newpage
} {
  Основные закономерности
} {
  \begin{center}
    \begin{tabular}{|c|c|c|c|c|c|} \hline
       & 1 & 2 & 3 & 4 & 5 \\ \hline
      $Q_1$ &  &  &  &  &  \\ \hline
      $Q_2$ &  &  &  &  &  \\ \hline
      $Q_3$ &  &  &  &  &  \\ \hline
    \end{tabular}
  \end{center}
  % \vspace{3cm}
} {
  \begin{center}
    \setlength{\tabcolsep}{10pt}\large
    \begin{tabular}{|c|c|c|c|c|} \hline
      $q_1$ & $q_2$ & $q_3$ & $q_4$ & $q_5$ \\ \hline
      & & & & \\ \hline
    \end{tabular}
  \end{center}
} {1}
\end{document}
